\section{Specialization}\label{sec:introduction}
\subsection[ ]{State of the Art and Limitations}

\begin{frame}{\secname}
	\framesubtitle{\subsecname}

	Systems programming aims to combine generic programming with type-specific \alert{optimizations}.
	% This concept is known as specialization and it allows developers to write flexible code that can be optimized for specific types or use cases.

	\pause
	\bigskip

	Mechanisms vary widely, from mature and widely used systems % e.g. C++ templates, which have been part of the language for decades
	to experimental features, % e.g. Rust or other languages, where specialization is still under active development.
	with varying degrees of success.

	\pause
	\bigskip

	In many modern ecosystems, specialization is often \alert{unstable}, forcing a tradeoff between performance and compiler stability.
	% The tradeoff is quite crucial, as leaving the stable channel may compromise the reliability of the software being developed and is also not always feasible in production environments.
	% On the other hand, the lack of specialization can lead to performance degradation, especially in performance-critical applications.

\end{frame}